\section{Numerical simulations}
\label{sec:rewrite-quantitative}

\subsection{Experimental design and initial conditions}
\label{subsec:rewrite-rsi-activation}

In this section, I compare transitional dynamics under four parameter configurations corresponding to the cases in Proposition~\ref{prop:rewrite-equilibrium-regimes}. I study an unanticipated activation of RSI at date zero. Before date zero, the economy already uses AI production services, but AI efficiency is fixed at $B_0$. For the chosen value of $B_0$, each economy starts on a labor-bottleneck balanced growth path. I compute this path with $\chi=0$, obtaining the corresponding initial capital stock $K_0$ for each parameter configuration. Setting $\chi=0$ rules out improvements in AI efficiency, not the expansion of AI services through inference compute. 

I compare four economies with $\sigma=0.9,1.0,1.1,1.5$, holding all other parameters and the initial values $A_0$, $N_0$, and $B_0$ fixed. For each economy, I choose $K_0$ to place it on its pre-RSI balanced growth path. In particular, for each elasticity, I choose the pre-event capital stock to satisfy
\begin{equation}
 r_0=\rho+\gamma,\qquad
 \frac{K_0}{Y_0}=\frac{\alpha}{\rho+\gamma+\delta}=3.30.
 \label{eq:rewrite-rsi-initial-ratio}
\end{equation}

Table~\ref{tab:rewrite-rsi-parameters} reports the parameter values and initial stocks used for the simulations. I use standard values for $\alpha$, $\delta$, and $\rho$, while $n$ and $\gamma$ are guided by long-run population and productivity projections. The AI production weight $\omega_X=0.10$  is an illustrative choices, not an empirical estimate. The research exponent $\eta=0.20$ is also illustrative but motivated by the inequality $\eta<\alpha$, which rules out the case studied in Proposition~\ref{prop:rewrite-research-scale}. The common upper bound on AI efficiency satisfies $\mathcal B(1.5)<\overline B<\mathcal B(1.1)$. Thus, AI and labor are substitutes in both the $\sigma=1.1$ and $\sigma=1.5$ cases, but only the latter converges to the AI-dominated regime.

I calibrate $\chi$ to match the AI research-expenditure
as a share of GDP, which is used as a proxy for $M/Y$. 

$0.182916\%$ have different origins. I originally calibrated
the central value, $\chi=7.616305$, to generate a $40\%$ fall in $p_X$ over
27 months: half the nearly $80\%$ decline in quality-adjusted AI-service
prices reported by \citet{oecdaimarkets2026} between January 2024 and
April 2026. Using half the observed decline is an illustrative choice.
I retain this value because it also yields a first-year research-expenditure
share of $0.182916\%$, which rounds to the central target of $0.183\%$.
This share is not an observed estimate or an independently matched moment.
For the slow transition, the reference is US training-compute expenditure
in 2023, estimated by \citet{korinekmckelvey2026} at $0.0664\%$ of
GDP.\footnote{The estimate is \$18.46 billion divided by
nominal GDP of \$27,811.517 billion from the BEA series \texttt{GDPA},
accessed through \href{https://fred.stlouisfed.org/data/GDPA}{FRED}
on September 14, 2026.}
The value $\chi=1.4378$ yields $0.06374\%$, an approximate match to an uncertain
proxy for research compute, not a direct measure of autonomous RSI.
Both research shares divide total research expenditure by total output
during the first year, rather than measuring $M_0/Y_0$.
I calibrate at $\sigma=1$ and hold $\chi$ fixed across the other three
elasticities within each exercise.

The following two subsections differ only in the calibration of $\chi$:
an illustrative research-expenditure target guides the central scenario, and a
lower research-expenditure target gives a slow-transition sensitivity.

The two values of $\chi$ have different origins. I originally calibrated
the central value, $\chi=7.616305$, to generate a $40\%$ fall in $p_X$ over
27 months: half the nearly $80\%$ decline in quality-adjusted AI-service
prices reported by \citet{oecdaimarkets2026} between January 2024 and
April 2026. Using half the observed decline is an illustrative choice.
I retain this value because it also yields a first-year research-expenditure
share of $0.182916\%$, which rounds to the central target of $0.183\%$.
This share is not an observed estimate or an independently matched moment.
For the slow transition, the reference is US training-compute expenditure
in 2023, estimated by \citet{korinekmckelvey2026} at $0.0664\%$ of
GDP.\footnote{The estimate is \$18.46 billion divided by
nominal GDP of \$27,811.517 billion from the BEA series \texttt{GDPA},
accessed through \href{https://fred.stlouisfed.org/data/GDPA}{FRED}
on September 14, 2026.}
The value $\chi=1.4378$ yields $0.06374\%$, an approximate match to an uncertain
proxy for research compute, not a direct measure of autonomous RSI.
Both research shares divide total research expenditure by total output
during the first year, rather than measuring $M_0/Y_0$.
I calibrate at $\sigma=1$ and hold $\chi$ fixed across the other three
elasticities within each exercise.

In Equation~\eqref{eq:rewrite-rsi-initial-ratio}, $Y_0$ is obtained from the
production and monopoly conditions with AI
already present. Thus $K_0$ differs across elasticities, while the initial
capital--output ratio and interest rate are common. With $A_0=N_0=1$ and
the common $B_0$, the four capital stocks are approximately $3.434$, $4.649$,
$5.037$, and $5.505$. These are calibrated alternative economies, not
different outcomes of introducing AI into one economy with the same capital
stock. The ratio is a theoretical BGP reference, not an exact empirical fit:
the PWT 11.0 comparison yields a 2023 current-price capital--output ratio
of about $3.61$ for the covered world sample and $3.42$ for the United States,
with capital definitions that differ from the one-good model \citep{pwt110}.

Before activation, consumption satisfies
\begin{equation}
 C(0^-)=Y_0-U_0-(\delta+n+\gamma)K_0>0.
 \label{eq:rewrite-rsi-pre-consumption}
\end{equation}
Only $K_0$ and $B_0$, together with the exogenous $A_0,N_0$, are inherited
by the post-event equilibrium. The BVP selects $C(0^+)$ and $q(0^+)$ anew.
Since the production block and its predetermined inputs do not change,
$Y$, $X$, $U$, $w$, $r$, and $p_X$ are continuous at activation. Research,
consumption, investment, and growth rates can adjust immediately. The
pre-event BGP requires an unanticipated activation: announcing RSI earlier
would generally change saving and consumption before date zero, as in the
technology-news mechanisms reviewed by \citet{beaudryportier2014}.

The sudden activation is an experimental assumption, not a requirement of
the RSI technology. It differs from the continued endogenous automation in
\citet{jonestonetti2026}. The exercise also omits the costs of expanding
compute capacity in GATE \citep{erdiletal2025gate} and the complementary
intangible investment studied by \citet{brynjolfssonrocksyverson2021}.
Those mechanisms can delay adjustment or measured productivity gains, but
they do not explain the paths computed here. Reducing $\chi$ does not remove
the unanticipated activation or require consumption and research spending
to be continuous at date zero. Appendix~\ref{app:rewrite-algorithm}
describes the equilibrium solver, numerical checks and replication files.

\begin{rewriteParameterTable}{Simulation parameters and initial stocks}{tab:rewrite-rsi-parameters}
$\omega_X$ & 0.10 & AI production weight & Illustrative.\\
$\omega_L$ & 0.90 & Effective-labor weight & Illustrative.\\
$\chi$ & 7.616305 (central);\newline 1.4378 (slow) & Research productivity after activation & First-year research-expenditure targets at $\sigma=1$: $0.183\%$ (central); approximately $0.065\%$ of GDP (slow transition).\\
$\sigma$ & 0.90; 1; 1.10; 1.50 & AI--labor substitution & Illustrative.\\
$\overline B$ & 1,361 & AI-efficiency upper bound & $1.10\mathcal B(1.50)$.\\
\midrule
$A_0,N_0$ & 1; 1 & Initial technology and population & Normalizations.\\
$B_0$ & 13.6 & Initial AI efficiency & $0.01\overline B$.\\
$K_0$ & 3.433645; 4.649320; 5.036956; 5.504667 & Initial capital & Fixed-efficiency BGP stocks with existing AI; Equation~\eqref{eq:rewrite-rsi-initial-ratio}.\\
$K_0/Y_0$ & 3.30 & Initial capital--output ratio & Derived from the pre-RSI BGP and preserved at activation.\\
\end{rewriteParameterTable}

